\documentclass[11pt]{article}
\usepackage[margin=1in]{geometry}
\usepackage{booktabs}
\usepackage{array}
\usepackage{xcolor}
\usepackage{hyperref}
\usepackage{amsmath}
\usepackage{enumitem}
\usepackage{titlesec}
\usepackage{longtable}
\usepackage{parskip}

\definecolor{linknavy}{HTML}{2a4d8f}
\hypersetup{colorlinks=true, linkcolor=linknavy, urlcolor=linknavy, citecolor=linknavy}

\titleformat{\section}{\Large\bfseries}{\thesection}{0.6em}{}
\titleformat{\subsection}{\large\bfseries}{\thesubsection}{0.6em}{}

\setlist[itemize]{topsep=2pt, itemsep=2pt, parsep=0pt}

\title{\texorpdfstring{\texttt{pi\_multik}}{pi\_multik}: Multi-Scale Persistence-Image Regression\\
\large Model Architectures, Data Pipeline, and Nested-Thomas Results}
\author{point-process-tda}
\date{\today}

\begin{document}
\maketitle
\vspace{-2em}

\tableofcontents

\section{Pipeline overview}

Every model in this report solves the same task: given a 2D point cloud sampled from a
stochastic point process, estimate the generative parameters that produced it (a
regression problem, not classification). The shared pipeline is:

\begin{enumerate}
\item \textbf{Sample} a point cloud from a process (Matérn hard-core, nested Thomas, \ldots)
  over the unit square.
\item \textbf{Filter} the cloud through a Distance-to-Measure (DTM) weighted-Rips
  filtration, at one or more density scales, producing a persistence diagram
  ($H_0$ + $H_1$ birth--death pairs) per scale.
\item \textbf{Vectorize} each diagram into a persistence image (PI): a fixed-size
  raster in (birth, persistence) coordinates, built by summing a Gaussian kernel at
  each off-diagonal point.
\item \textbf{Regress} the process parameters from the stacked per-scale PIs (plus a
  couple of scalar side-features) with a CNN, trained with MSE loss in
  standardized (log-z-score) target space.
\end{enumerate}

The \texttt{pi\_multik} family (Section~\ref{sec:architectures}) is this repo's
name for the family of architectures that fuse \emph{several DTM scales at once}
per cloud, as opposed to using a single scale in isolation.

\section{Model architectures}
\label{sec:architectures}

All variants share the same encoder building block and the same MLP regression
head; they differ only in \emph{how the per-scale embeddings are fused}.

\subsection{Shared building blocks}

Exact sizes below are what every config in \texttt{configs/runs/} actually
uses (matérn and nested-Thomas configs are identical on every value here) — not
just code defaults.

\begin{itemize}
  \item \textbf{CoordConv PI encoder} (\texttt{CoordConvPIEncoder}): input is one
    $(H_0, H_1)$ image pair — \texttt{in\_channels}$=2$, resolution $64\times64$ —
    for a single DTM scale. Two constant coordinate channels ($x,y \in [-1,1]$)
    are concatenated first, bringing the first \texttt{Conv2d}'s input to
    $2+2=4$ channels. Three \texttt{Conv2d(k=3)--ReLU} blocks with
    \texttt{conv\_channels}$=(32, 64, 128)$ (i.e.\ $4\to32\to64\to128$ channels),
    \texttt{MaxPool2d}+\texttt{Dropout2d}($p=0.2$) between them, then global
    average pooling ($128\times64\times64 \to 128$) and a final
    \texttt{Linear(128, 64)} down to a fixed \texttt{embedding\_dim}$=64$ vector
    — one 64-d embedding per scale.
  \item \textbf{Regression head} (\texttt{ParameterEstimator}): a plain MLP —
    \texttt{Linear--ReLU--Dropout($p=0.1$)} $\times 2$ with
    \texttt{head\_hidden\_dims}$=(64, 32)$, then a final linear layer to the
    target dimension ($2$ for Matérn, $5$ for nested Thomas). Identical across
    every variant below; only its \emph{input} width changes (see each
    variant's head-input size below).
  \item \textbf{Extra scalar side-channel}: $\log N$ (point count) always,
    giving \texttt{n\_extra}$=1$ in every run in this report
    (\texttt{include\_entropy: false} throughout); optionally, per-(dimension,
    scale) persistence entropy would add more columns here if turned on.
    Concatenated onto the fused embedding just before the head.
\end{itemize}

\subsection{\texttt{pi\_multik} --- flat-concatenation (late-fusion) baseline}
\begin{itemize}
  \item One \emph{shared-weight} \texttt{CoordConvPIEncoder} instance
    (\texttt{embedding\_dim}$=64$, as above), applied to every scale's
    $(H_0,H_1)$ pair independently (scales folded into the batch dimension, so
    it is one conv-stack call per forward pass, not $K$ separate ones).
  \item The $K$ resulting 64-d embeddings are concatenated flat into a single
    $K\times64$ vector — order-blind fusion, no notion of scale adjacency. Head
    input size is therefore $K\cdot64 + 1$: e.g.\ $193$ for the best nested-Thomas
    model ($K=3$, $k\in\{5,10,15\}$), or $449$ for the full 7-scale union.
  \item That $K\cdot64+1$ vector is what the shared head (above) actually
    consumes.
  \item \textbf{\texttt{pi\_multik\_earlyfusion}} is the sibling that instead
    stacks every scale's channels together \emph{before} the first conv layer —
    one \texttt{CoordConvPIEncoder} with \texttt{in\_channels}$=2K+2$ (e.g.\
    $8$ for $K=3$, $16$ for $K=7$, including the $+2$ coordinate channels),
    producing a single fixed $64$-d embedding regardless of $K$ (head input
    $64+1=65$, \emph{not} $K$-dependent, unlike the late-fusion baseline above).
    This was the original \texttt{pi\_multik} design; it is registered and
    actively trained today as its own sibling method (Section~\ref{sec:results}
    reports real 10-seed results for it, not just an architectural aside).
\end{itemize}

\subsection{\texttt{pi\_multik\_scaleconv} --- 1D conv fusion over the scale axis}
\begin{itemize}
  \item Identical data path and shared-weight encoder as \texttt{pi\_multik}
    (same \texttt{embedding\_dim}$=64$ per-scale embeddings).
  \item Instead of flat concatenation, the $K$ 64-d embeddings are treated as
    an \emph{ordered sequence} $(K, 64)$ along the scale axis and passed
    through \texttt{ScaleConvFusion}: \texttt{Conv1d}$(64\to\texttt{scale\_fusion\_hidden}{=}128,
    k{=}3,\text{pad}{=}1)$, LayerNorm, ReLU, then
    \texttt{Conv1d}$(128\to\texttt{scale\_fusion\_out\_dim}{=}128, k{=}3,\text{pad}{=}1)$,
    LayerNorm, ReLU — the padding keeps the scale axis at length $K$ throughout —
    then \texttt{AdaptiveAvgPool1d(1)} collapses it, giving a single fixed
    $128$-d vector \emph{regardless of $K$}. Head input is therefore always
    $128+1=129$, whether $K=3$ or $K=7$ — the only one of the four architectures
    whose head-input size doesn't grow with the number of scales fused.
  \item This is the only variant that both mixes information \emph{across}
    adjacent scales (via the \texttt{Conv1d}) \emph{and} pools the scale axis
    away afterward — \texttt{pi\_multik\_towers} (next) also mixes adjacent
    scales the same way, but does not pool.
\end{itemize}

\subsection{\texttt{pi\_multik\_towers} --- independent per-scale encoders + tower fusion}
\begin{itemize}
  \item Unlike the two variants above, each scale gets its \emph{own}
    (non-shared-weight) \texttt{CoordConvPIEncoder} — $K$ independent "towers",
    each still producing a $64$-d embedding, but from $K$ separate sets of
    conv weights rather than one shared set.
  \item Embeddings are stacked into a $(K, 64)$ sequence and passed through
    \texttt{TowerConv} — the identical \texttt{Conv1d}$(64\to128\to128, k{=}3)$
    + LayerNorm + ReLU recipe as \texttt{ScaleConvFusion} (own weights, not
    shared with it), plus a \texttt{Dropout1d} — but \textbf{without} the final
    pooling step: the $(128, K)$ output is flattened, not averaged, to
    $\texttt{out\_dim}\times K = 128K$. Head input is therefore $128K+1$: $385$
    for $K=3$ ($k\in\{5,10,15\}$), $897$ for $K=7$ — growing with $K$ exactly
    the way \texttt{pi\_multik}'s does, just with a $128$-wide block per scale
    instead of $64$.
  \item Strictly more parameters than the other three variants: $K$ independent
    encoders (rather than one shared or one merged-input encoder) \emph{and}
    the largest head-input width of the four for $K>1$.
\end{itemize}

\subsection{\texttt{pi\_multik\_fusion} --- TDA + classical-statistics fusion}
\begin{itemize}
  \item Adds a second branch replicating Vihrs (2022)'s own architecture: a
    3-layer \texttt{Conv1d}$(1\to64\to64\to64, k{=}7)$ stack (with two
    \texttt{MaxPool1d(5)}s) over the $513$-point centred Ripley $L(r)-r$
    curve, flattened and passed through a \texttt{Linear} layer down to the
    \emph{same} \texttt{embedding\_dim}$=64$ used everywhere else in this
    report, so it can be concatenated on equal footing with the image branch.
  \item Image branch is the plain \texttt{pi\_multik} flat-concat path verbatim
    (shared-weight \texttt{CoordConvPIEncoder}, $K\times64$ output).
  \item All three pieces concatenate into one head input:
    $\underbrace{64}_{L(r)-r\text{ branch}} + \underbrace{K\cdot64}_{\text{image branch}} + \underbrace{1}_{\log N}$
    — e.g.\ $64+192+1=257$ for $K=3$.
  \item Implemented and registered, but \textbf{not yet trained on the current
    Matérn / nested-Thomas datasets} (no \texttt{results/} for it under either
    process) — included here for architectural completeness only.
\end{itemize}

\subsection{\texttt{vihrs} --- classical-statistics baseline (not a \texttt{pi\_multik} variant)}
\begin{itemize}
  \item Independent re-implementation of Vihrs (2022): a single-scale
    \texttt{Conv1d} CNN over the isotropic-edge-corrected $L(r)-r$ curve
    (513 points, $r_{\max}=0.25$) concatenated with $\log N$, no persistence
    diagrams or TDA features at all.
  \item Used throughout this report purely as an external reference point — "does
    the TDA machinery beat a much simpler classical-statistics feature set?"
\end{itemize}

\section{How the data is used}

\begin{itemize}
  \item \textbf{Input cloud}: 2D points in $[0,1]^2$, $N$ ranging from tens to
    roughly a thousand depending on the process and sampled parameters.
  \item \textbf{DTM scale}: parameterized by a \emph{mass fraction} $m \in (0,1]$
    rather than a fixed neighbor count $k$, because a fixed $k$ means a very
    different fraction of each cloud once $N$ varies across the parameter sweep.
    Per cloud, $k = \mathrm{round}(m \cdot N)$ is recomputed from its own $N$, so
    the same $m$ means "the same relative neighborhood size" for every cloud
    regardless of point count. The swept grid used throughout this report is
    $m \in \{0.01, 0.02, 0.04, 0.10, 0.20, 0.45, 0.90\}$ — fine (small, local,
    noise-sensitive neighborhoods) to coarse (large, smoothed, near-parent-scale
    neighborhoods).
  \item \textbf{Persistence image}: $64\times64$ raster per (scale, homology
    dimension); axis bounds (birth/persistence range) are calibrated once from
    percentiles of the \emph{training} diagrams' own distribution (95th/99.9th),
    not hand-picked. Each pixel is the exact integral of the diagram's
    persistence surface (a 2-D Gaussian per off-diagonal point, weighted by
    \texttt{weight\_fn}) over that pixel's box, via the Gaussian CDF — not a
    point sample at the pixel center. Gaussian bandwidth is a single fixed
    \texttt{sigma\_pixels}$=2.0$, shared by the birth and persistence axes and
    shared across every diagram (an earlier per-diagram \emph{adaptive}
    $\sigma$ design was retired in favor of this; see
    Section~\ref{sec:sigma}).
  \item \textbf{Channel stacking across scales}: per-scale PI pairs are aligned by
    cloud seed and intersected (a diagram can be empty and get filtered out at one
    scale but not another) before being stacked into the $(N, K, 2, 64, 64)$ tensor
    the network consumes.
  \item \textbf{Targets}: the process's own generative parameters — 2 for Matérn
    (\texttt{parent\_intensity}, \texttt{hardcore\_radius}), 5 for nested Thomas
    (\texttt{parent\_intensity} [meta-layer density], \texttt{meta\_offspring},
    \texttt{meta\_cluster\_scale}, \texttt{mean\_offspring}, \texttt{cluster\_scale}).
    All are log-transformed then z-scored, with the transform fit on the training
    split only and frozen for val/test/adversarial.
  \item \textbf{Splits}: 70/15/15 train/val/test (seeded per run), plus a separate
    \emph{frozen adversarial} 12.5\% hold-out whose underlying parameter vectors
    never appear in train/test at all (not just different noise realizations) —
    it tests generalization to unseen parameter combinations, not memorization.
\end{itemize}

\section{Key hyperparameters}

\begin{center}
\begin{tabular}{@{}ll@{}}
\toprule
\textbf{Hyperparameter} & \textbf{Role / value used here} \\
\midrule
DTM scale grid ($m$ values) & the single most consequential knob — see Section~\ref{sec:results} \\
Homology dimensions & $H_0$, $H_1$ (both used everywhere in this report) \\
PI resolution & $64 \times 64$ \\
\texttt{sigma\_pixels} & $2.0$ — fixed Gaussian kernel bandwidth, in pixels,
shared by the birth and persistence axes (retired the earlier per-diagram
adaptive design; see Section~\ref{sec:sigma}) \\
\texttt{embedding\_dim} & 64 (per-scale CoordConv embedding size) \\
\texttt{conv\_channels} & $(32, 64, 128)$ \\
\texttt{scale\_fusion\_hidden} / \texttt{out\_dim} & 128 / 128 (scaleconv, towers) \\
\texttt{scale\_fusion\_kernel\_size} & 3 \\
Head \texttt{hidden\_dims} / dropout & $(64, 32)$ / $0.1$ \\
Batch size & 32 \\
Optimizer & Adam, $\text{lr}=10^{-3}$, weight decay $10^{-4}$ \\
Max epochs / early-stop patience & 500 / 30 (no val-loss improvement) \\
Seeds & up to 10 per configuration (see per-table $n$) \\
Train/val/test split & 70/15/15 \\
Adversarial hold-out & 12.5\% of parameter vectors, disjoint from train/test \\
\bottomrule
\end{tabular}
\end{center}

\section{Results: Nested Thomas vs.\ \texttt{vihrs}}
\label{sec:results}

Loss is MSE in standardized (log-z-score) target space, averaged over the five
nested-Thomas targets — $1.0$ is roughly what predicting the training mean
would score, so lower is better and values noticeably below 1 indicate the
model is extracting real signal. \textbf{Matérn results are omitted from this
revision of the report}: nested Thomas below has been fully recomputed with
the fixed persistence imager and retrained end to end, while Matérn's own
rerun on the same fixed imager is still pending — it will be reinstated once
that rerun completes.

\subsection{Channel-combination landscape}
\label{sec:landscape}

Dozens of channel (DTM-scale) combinations were tried; this is not all of them,
just enough to see the shape of the landscape. Every row below uses the base
\texttt{pi\_multik} architecture, so the comparison isolates \emph{which
scale(s) feed the network} from \emph{which fusion architecture} combines them
(that second axis is examined separately in Section~\ref{sec:bestvsvihrs}).

\paragraph{Single DTM scale.}
\begin{center}
\begin{tabular}{@{}lrr@{}}
\toprule
\textbf{DTM scale $m$} & \textbf{test loss} & \textbf{adv.\ loss} \\
\midrule
0.01 & 0.5245 $\pm$ 0.0091 & 0.5157 $\pm$ 0.0063 \\
0.02 & 0.5173 $\pm$ 0.0083 & 0.5017 $\pm$ 0.0033 \\
0.04 & 0.5287 $\pm$ 0.0067 & 0.5113 $\pm$ 0.0018 \\
0.10 & 0.5500 $\pm$ 0.0089 & 0.5345 $\pm$ 0.0039 \\
0.20 & 0.5688 $\pm$ 0.0103 & 0.5584 $\pm$ 0.0044 \\
0.45 & 0.5788 $\pm$ 0.0141 & 0.5710 $\pm$ 0.0102 \\
0.90 & 0.5990 $\pm$ 0.0108 & 0.5969 $\pm$ 0.0151 \\
\bottomrule
\end{tabular}
\end{center}
\emph{$n=3$ seeds per cell; method \texttt{pi\_multik}.}

\paragraph{Multi-scale unions, and the alternative fixed-$k$ sweep.}
\begin{center}
\begin{tabular}{@{}lrr@{}}
\toprule
\textbf{Channel set} & \textbf{$n$} & \textbf{Nested Thomas test loss} \\
\midrule
fine union $m\in\{0.01,0.02,0.04\}$              & 3  & 0.5077 $\pm$ 0.0075 \\
spread union $m\in\{0.01,0.10,0.90\}$            & 3  & 0.5125 $\pm$ 0.0088 \\
coarse union $m\in\{0.20,0.45,0.90\}$            & 3  & 0.5641 $\pm$ 0.0157 \\
5-scale union $m\in\{0.01,0.04,0.10,0.20,0.45\}$ & 3  & 0.5059 $\pm$ 0.0104 \\
full 7-scale union (all swept $m$)               & 3  & 0.5102 $\pm$ 0.0083 \\
\midrule
\textbf{fixed-$k$ sweep, $k\in\{5,10,15\}$}       & \textbf{10} & \textbf{0.5014 $\pm$ 0.0082} \\
\bottomrule
\end{tabular}
\end{center}
\emph{All rows: method \texttt{pi\_multik}. $k\in\{5,10,15\}$ is this repo's
original, pre-mass-fraction DTM parameterization (fixed neighbor count rather
than $m=k/N$ recomputed per cloud) — included here since it predates, and is
directly comparable against, the $m$-based sweep.}

\textbf{Reading this landscape:} nested Thomas degrades steadily as the DTM
scale $m$ grows, from a best single scale of $m=0.02$ ($0.5173$) down to
$m=0.90$'s $0.5990$ — finer neighborhoods carry more signal than coarser ones,
consistently. \emph{Fusing multiple scales together does not clearly beat
just using the single best fine scale}: every union that includes at least
one fine scale lands in a tight $0.506$--$0.513$ band regardless of exactly
which scales are combined, while the one union built entirely from coarse
scales ($m\in\{0.20,0.45,0.90\}$) is clearly worse ($0.5641$) — confirming the
fine scales are carrying essentially all of the union's signal. The
fixed-$k$ sweep ($k\in\{5,10,15\}$) is the best configuration in the whole
landscape ($0.5014$), a clear improvement (not just a tie, as in the previous
persistence-image version of this comparison) over the next-best 5-scale
union ($0.5059$). It is also the most heavily seeded ($n=10$) configuration,
which is why it is used as \emph{the} best model in
Section~\ref{sec:bestvsvihrs}.

\subsection{Best model vs.\ \texttt{vihrs}}
\label{sec:bestvsvihrs}

Best-model channel set: the fixed-$k$ sweep ($k\in\{5,10,15\}$,
Section~\ref{sec:landscape}), all four \texttt{pi\_multik} architectures at
$n=10$ — matched exactly to \texttt{vihrs}'s own $n=10$ (\texttt{vihrs} itself
does not depend on persistence images, so its numbers are unchanged by this
recompute):

\begin{center}
\begin{tabular}{@{}lrrr@{}}
\toprule
\textbf{Method} & \textbf{$n$} & \textbf{test loss} & \textbf{adv.\ loss} \\
\midrule
\texttt{pi\_multik}              & 10 & \textbf{0.5014 $\pm$ 0.0082} & \textbf{0.4875 $\pm$ 0.0025} \\
\texttt{pi\_multik\_earlyfusion} & 10 & 0.5015 $\pm$ 0.0060 & 0.4879 $\pm$ 0.0029 \\
\texttt{pi\_multik\_scaleconv}   & 10 & 0.5035 $\pm$ 0.0066 & 0.4902 $\pm$ 0.0016 \\
\texttt{pi\_multik\_towers}      & 10 & 0.5076 $\pm$ 0.0085 & 0.4928 $\pm$ 0.0029 \\
\texttt{vihrs}                   & 10 & 0.5224 $\pm$ 0.0108 & 0.5149 $\pm$ 0.0034 \\
\bottomrule
\end{tabular}
\end{center}

\textbf{Every \texttt{pi\_multik} architecture clearly beats \texttt{vihrs}}
— even the weakest variant (\texttt{towers}, $0.5076$) sits comfortably below
\texttt{vihrs}'s $0.5224$, and the gap between the best TDA model and
\texttt{vihrs} ($\approx 0.021$) is bigger than either side's standard
deviation. Plain \texttt{pi\_multik} (flat-concat late fusion) and
\texttt{pi\_multik\_earlyfusion} are now essentially tied for best
($0.5014$ vs.\ $0.5015$, well within one standard deviation of each other) —
the fancier fusion mechanisms (scale-aware convolution, independent per-scale
towers) still don't improve on the simplest architecture here.

\subsection{Per-target breakdown (best model vs.\ \texttt{vihrs})}

\begin{center}
\begin{tabular}{@{}lrr@{}}
\toprule
\textbf{Target} & \textbf{best \texttt{pi\_multik} model} & \texttt{vihrs} \\
\midrule
\texttt{parent\_intensity} (meta-layer density) & 0.8740 & 0.8765 \\
\texttt{meta\_offspring}      & 0.8546 & 0.8600 \\
\texttt{meta\_cluster\_scale} & 0.5261 & 0.5607 \\
\texttt{mean\_offspring}      & 0.2109 & 0.2754 \\
\texttt{cluster\_scale}       & 0.0413 & 0.0395 \\
\bottomrule
\end{tabular}
\end{center}
\emph{Best model: \texttt{pi\_multik}, $k\in\{5,10,15\}$; $n=10$ for both
columns.}

\textbf{Notable pattern:} on the two \emph{meta}-layer / coarse parameters
(\texttt{parent\_intensity}, \texttt{meta\_offspring}), \texttt{pi\_multik} and
\texttt{vihrs} remain statistically indistinguishable ($0.85$--$0.88$ for both,
on both targets) — neither the TDA feature set nor the classical $L(r)-r$
feature set is extracting much signal about the coarse scale of this
two-level process. On the \emph{fine}-layer / geometric parameters the
picture is mixed: \texttt{pi\_multik} is clearly ahead on
\texttt{meta\_cluster\_scale} ($0.53$ vs.\ $0.56$) and, most visibly, on
\texttt{mean\_offspring} ($0.21$ vs.\ $0.28$); but on \texttt{cluster\_scale}
the two are now statistically tied ($0.0413$ vs.\ $0.0395$, well within one
standard deviation of each other) rather than \texttt{pi\_multik} holding a
clear edge as in the previous persistence-image version of this report. Put
together: \texttt{pi\_multik}'s overall advantage over the classical baseline
(Section~\ref{sec:bestvsvihrs}) is concentrated in
\texttt{meta\_cluster\_scale} and \texttt{mean\_offspring} specifically, not a
blanket improvement across every parameter — and the shared difficulty on the
two coarse parameters looks more like a ceiling in the \emph{data or process
design} (or in every feature set's ability to see the meta-scale at all, at
the window sizes tried so far) than a model-architecture problem, since it
shows up identically in two very different feature representations.

\subsection{Note on the persistence-image sigma: adaptive retired in favor of a fixed value}
\label{sec:sigma}

The per-diagram \emph{adaptive} $\sigma$ used in earlier versions of this
report (\texttt{sigma\_frac}$\times$the diagram's own median off-diagonal
persistence, floored at \texttt{min\_sigma\_pixels}) has been retired. It is
still available as \texttt{adaptive\_persistence\_image.py} for comparison,
but is no longer what any config in \texttt{configs/runs/} uses. Two problems
motivated dropping it, both checked on nested Thomas, $m=0.10$
(\texttt{notebooks/pi\_sigma\_diagnostics.ipynb}):

\begin{itemize}
  \item \textbf{Not meaningfully adaptive.} At resolution 64 and
    \texttt{sigma\_frac}$=0.1$, the 1-pixel floor bound for essentially
    \emph{every} $H_0$ diagram and $\approx$99.7\% of $H_1$ diagrams — i.e.\
    $\sigma$ sat at its floor almost everywhere, engaging only on the long
    tail of unusually high-persistence-spread diagrams.
  \item \textbf{Weaker empirical stability.} Under cloud-jitter perturbation,
    $\|\Delta\text{PI}\|_2$ correlated far more weakly with
    bottleneck/Wasserstein distance for the adaptive design (Pearson
    $r\approx0.45$--$0.51$ for $H_0$) than for a fixed-$\sigma$ comparator
    ($r\approx0.90$--$0.96$) — evidence the diagram-dependent $\sigma$ was
    degrading the Lipschitz stability property persistence images are
    supposed to have (Adams et al., 2017), not just under-engaging.
\end{itemize}

The replacement, \texttt{PersistenceImager} in \texttt{persistence\_image.py},
uses a single \texttt{sigma\_pixels}$=2.0$ shared by both axes and every
diagram, plus exact CDF-based box integration per pixel rather than a point
sample at the pixel center. The value $2.0$ was chosen from a
speckle/connected-component diagnostic sweep (\texttt{sigma\_pixels}
$\in[0.25,6]$) over the real \texttt{dtm\_k5/10/15} diagrams: $H_1$ stays
heavily speckled (top-1\%-pixel-energy fraction $>0.9$) below
\texttt{sigma\_pixels}$=2$, while $H_0$ starts collapsing toward a single
connected component for a growing share of diagrams above
\texttt{sigma\_pixels}$=2$--$3$ — $2.0$ sits at that balance point (see
\texttt{figures/pi\_adaptive\_sigma\_comparison.png}, comparing the retired
adaptive design against this fixed value). \textbf{Caveat: this choice has
not yet been cross-checked against training/test loss} —
\texttt{scripts/sweep\_pi\_sigma.py} exists for exactly that sweep but has not
been run yet. All results in Section~\ref{sec:results} use
\texttt{sigma\_pixels}$=2.0$ throughout; re-running that sweep is the natural
next step before treating $2.0$ as settled.

\section{Figures prepared for the presentation}

Four standalone PNGs are in \texttt{docs/figures/}, generated directly from this
repo's own data/pipeline code (not mockups):

\begin{itemize}
  \item \textbf{\texttt{nested\_thomas\_example.png}} --- one realized nested
    Thomas cloud (an actual dataset cloud, not synthetic), points colored by
    meta-cluster, with a zoomed inset on one meta-cluster showing its fine
    sub-clusters.
  \item \textbf{\texttt{persistence\_diagram\_example.png}} --- the $H_0$/$H_1$
    DTM persistence diagram ($m=0.10$) for that same cloud.
  \item \textbf{\texttt{persistence\_images\_multiscale.png}} --- that same
    cloud's persistence images across all 7 swept DTM scales, $H_0$ and $H_1$,
    showing the fine$\to$coarse progression.
  \item \textbf{\texttt{pi\_adaptive\_sigma\_comparison.png}} --- at one fixed
    DTM scale, fixed vs.\ adaptive $\sigma$ side by side on three clouds spanning
    the adaptive-$\sigma$ range, illustrating Section~\ref{sec:sigma}'s finding
    directly.
\end{itemize}

\end{document}
